\documentclass[10pt,a4paper]{amsart}
\usepackage[margin=19mm]{geometry}
\usepackage{amsmath,amssymb,mathtools}
\usepackage[scr=rsfs,cal=euler]{mathalfa}
\usepackage{xfrac}
\usepackage[unicode,hidelinks,bookmarks=false]{hyperref}
\newcommand{\ie}{i.e.\ }
\newcommand{\cf}{cf.\ }
\newcommand{\resp}{respectively\ }
\newcommand{\Subsection}{\subsection}
\providecommand{\nobreakdash}{\nobreak\hbox{-}}
\setlength{\emergencystretch}{2em}
\multlinegap0pt
\usepackage{amssymb}
\usepackage{enumerate}
\newtheorem{conjecture}{Conjecture}
\newtheorem{theorem}{Theorem}
\newtheorem{lemma}{Lemma}
\title{On a conjecture on Romanoff type sumsets}
\author{Yuchen Ding, Liangxun Li}
\date{Source: arXiv:2606.06118v2 (CC BY 4.0)}
\begin{document}
\maketitle

\noindent\emph{Source and licence.} On a conjecture on Romanoff type sumsets. Version arXiv:2606.06118v2 (CC BY 4.0), distributed by arXiv under the Creative Commons Attribution 4.0 International licence, http://creativecommons.org/licenses/by/4.0/, as recorded in the arXiv licence metadata for this exact version and shown on the abstract page https://arxiv.org/abs/2606.06118v2. Full licence notice: \texttt{licenses/arxiv-2606.06118-CC-BY-4.0.txt}.\par\smallskip

\noindent\textit{Editorial scope.} Counted unit: the article's theorem theorem-1 (source0.tex lines 200--207). The article's complete original proof of it is delivered with it (source0.tex lines 280--326, one \texttt{proof} environment of 47 lines, which the article places away from the statement). No mathematics is written, repaired or re-derived: the statement and the proof are the article's own lines, in the article's own notation, and no retained line is substituted or re-flowed.  Retained uncounted as the source-local context the proof needs: lines 177--183; lines 232--238; lines 253--259.  Nothing was omitted from any retained span, so the package carries no omission record.  No retained line contains a citation, so the wrapper carries no bibliography and no bibliography entry.  No retained line contains a figure environment, an image-inclusion command or drawing code, so no image is delivered and the package declares no asset dependency.  Editorial formatting only, in addition: an AMS \texttt{amsart} preamble that restates the article's own declarations and macros used by the retained lines, namely the environments \texttt{conjecture}, \texttt{theorem}, \texttt{lemma} on separate counters, together with the standard packages those lines need.  The retained environments print in this document as Conjecture 1, Theorem 1, Lemma 1 in order of appearance, whereas the article prints the same items with the numbering of its own full document.  The complete native source of the article as distributed in the arXiv e-print for this exact version is bundled unchanged as \texttt{sources/202160/source0.tex}.
\begin{conjecture}[Weak uniform Hardy-Littlewood Conjecture]\label{conjecture-2}
There is a fixed $\delta>0$ such that the inequality
$$
\#\big\{p\le x: p, p+h\in \mathbb{P}\big\}\gg_\delta \frac{x}{(\log x)^2}
$$
holds uniformly for any even $h\le x^\delta$, provided that $x$ is sufficiently large.
\end{conjecture}

\begin{theorem}\label{propo:2}
For any fixed positive integer $k$, we have
$$
\sum_{n\le x}f_{\mathcal{A}}(n)^k\ll x (\log x)^{k\left(\frac{1}{r_1}+\cdots+\frac{1}{r_t}-1\right)},
$$
where the implied constant depends at most on $\mathcal{A}$ and $k$.
\end{theorem}

\begin{lemma}\label{lem1}
Let $r_1,\cdots, r_t$ be positive number. Then for sufficiently large $x$ we have
$$
\mathcal{A}(x)\gg (\log x)^{\frac{1}{r_1}+\cdots+\frac{1}{r_t}},
$$
where the implied constant depends only on $r_1,\cdots, r_t$.
\end{lemma}

\begin{theorem}\label{theorem-1}
Let $r_1,\cdots, r_t$ are positive numbers with $r_i\ge 1~(1\le i\le t)$ and 
$$
r_1^{-1}+\cdots+r_t^{-1}\ge 1.
$$
 Then
 the set $\mathcal{R}_{\mathcal{A}}\cap(\mathcal{R}_{\mathcal{A}}+2)$ has a positive lower asymptotic density under the assumption of Conjecture \ref{conjecture-2}.
\end{theorem}

\begin{proof}[Proof of Theorem \ref{theorem-1}]
Let $\mathcal{R}_{\mathcal{A},2}=\mathcal{R}_{\mathcal{A}}\cap(\mathcal{R}_{\mathcal{A}}+2)$ for simplification. Using twice the Cauchy-Schwarz inequality we have
\begin{align}\label{eq-2-1}
\Big(\sum_{n\le x} f_{\mathcal{A}}(n)f_{\mathcal{A}}(n+2)\Big)^2&\le \Big(\sum_{\substack{n\le x\\ f_{\mathcal{A}}(n)f_{\mathcal{A}}(n+2)\ge 1}} 1^2 \Big)\Big(\sum_{n\le x} f_{\mathcal{A}}(n)^2f_{\mathcal{A}}(n+2)^2\Big)\nonumber\\
&\le \mathcal{R}_{\mathcal{A},2}(x)\sqrt{\sum_{n\le x} f_{\mathcal{A}}(n)^4}\sqrt{\sum_{n\le x} f_{\mathcal{A}}(n+2)^4}.
\end{align}
By Theorem \ref{propo:2}, we know
\begin{align}\label{eq-2-2}
\sum_{n\le x} f_{\mathcal{A}}(n)^4\ll x (\log x)^{4\left(\frac{1}{r_1}+\cdots+\frac{1}{r_t}-1\right)}
\end{align}
and
\begin{align}\label{eq-2-3}
\sum_{n\le x} f_{\mathcal{A}}(n+2)^4\le \sum_{n\le 2x} f_{\mathcal{A}}(n)^4 \ll x (\log x)^{4\left(\frac{1}{r_1}+\cdots+\frac{1}{r_t}-1\right)}.
\end{align}
Inserting  (\ref{eq-2-2}) and (\ref{eq-2-3}) into (\ref{eq-2-1}), we get
$$
\mathcal{R}_{\mathcal{A},2}(x)\gg \frac{\Big(\sum_{n\le x} f_{\mathcal{A}}(n)f_{\mathcal{A}}(n+2)\Big)^2}{x(\log x)^{4\left(\frac{1}{r_1}+\cdots+\frac{1}{r_t}-1\right)}}.
$$
Thus, to prove our theorem it suffices to show that
\begin{align}\label{eq-2-4}
\sum_{n\le x} f_{\mathcal{A}}(n)f_{\mathcal{A}}(n+2)\gg x(\log x)^{2\left(\frac{1}{r_1}+\cdots+\frac{1}{r_t}-1\right)}.
\end{align}

From now on, the symbols $p$'s will always be primes. Clearly,
\begin{align*}
\sum_{n\le x} f_{\mathcal{A}}(n)f_{\mathcal{A}}(n+2)&=\sum_{n\le x}\sum_{\substack{n=p_1+a_1\\ n+2=p_2+a_2\\ a_1, a_2\in \mathcal{A}}}1\\
&=\sum_{a_1, a_2\in \mathcal{A}}\sum_{p_1+a_1=p_2+a_2-2\le x}1\\
&\ge \sum_{\substack{a_1, a_2\in \mathcal{A}\\a_2<a_1\le x^\delta}}\sum_{\substack{p_1\le \frac{x}{2}\\ p_2=p_1+a_1-a_2+2}}1
\end{align*}
where $\delta$ is the one occurred in Conjecture \ref{conjecture-2}. By Conjecture \ref{conjecture-2}, we have
\begin{align*}
\sum_{\substack{p_1\le \frac{x}{2}\\ p_2=p_1+a_1-a_2+2}}1\gg \frac{x}{(\log x)^2}.
\end{align*}
Therefore, we conclude that
\begin{align*}
\sum_{n\le x} f_{\mathcal{A}}(n)f_{\mathcal{A}}(n+2)\gg \frac{x}{(\log x)^2}\sum_{\substack{a_1, a_2\in \mathcal{A}\\a_2<a_1\le x^\delta}}1=\frac{x}{(\log x)^2}\binom{\mathcal{A}\left(x^\delta\right)}{2}.
\end{align*}
By Lemma \ref{lem1} we know that
$$
\mathcal{A}\left(x^\delta\right)\gg (\log x)^{\frac{1}{r_1}+\cdots+\frac{1}{r_t}},
$$
from which we deduce that
$$
\sum_{n\le x} f_{\mathcal{A}}(n)f_{\mathcal{A}}(n+2)\gg x(\log x)^{2\left(\frac{1}{r_1}+\cdots+\frac{1}{r_t}-1\right)},
$$
proving (\ref{eq-2-4}).
\end{proof}

\end{document}
